\documentclass[11pt]{article}

\usepackage[margin=1in]{geometry}
\usepackage{amsmath,amssymb}
\usepackage{bm}
\usepackage{booktabs}
\usepackage{siunitx}
\usepackage[hidelinks]{hyperref}

\newcommand{\mjorbit}{mjorbit}
\newcommand{\mj}{MuJoCo}
\newcommand{\R}{\mathbb{R}}

\begin{document}

\section{Methods}
\label{sec:methods}

\mjorbit{} couples an Earth-orbit propagator to \mj{}'s articulated rigid-body
dynamics. The chief spacecraft state $(\bm R,\bm V)$ is stored in ECI
coordinates with km and \si{km/s} units. \mj{} itself is run in SI units in a
chief-centered world frame whose axes remain parallel to ECI. Thus, for body
$i$ with \mj{} world position and velocity $(\bm r_i,\bm v_i)$,
\begin{equation}
  \bm R_i = \bm R + 10^{-3}\bm r_i,\qquad
  \bm V_i = \bm V + 10^{-3}\bm v_i,
  \label{eq:absolute-state}
\end{equation}
are the body's absolute ECI position and velocity. Because the \mj{} world
frame translates with the chief but does not rotate, the only inertial
correction required in the multibody dynamics is the translational
pseudo-force associated with the chief acceleration.

\subsection{Reference Orbit and Gravity}
\label{subsec:reference-orbit}

The chief is propagated as a point mass under Earth's gravity and a
swarm-mean non-gravity acceleration,
\begin{equation}
  \dot{\bm R}=\bm V,\qquad
  \dot{\bm V}=\bm g(\bm R)+\bm a_{\rm fb}.
  \label{eq:chief-dynamics}
\end{equation}
The gravity model is two-body gravity with an optional $J_2$ correction,
\begin{equation}
  \bm g(\bm R)=
  -\frac{\mu}{r^3}\bm R
  +\frac{3\mu J_2R_\oplus^2}{2r^5}
  \begin{bmatrix}
    R_x(5R_z^2/r^2-1)\\
    R_y(5R_z^2/r^2-1)\\
    R_z(5R_z^2/r^2-3)
  \end{bmatrix},
  \qquad r=\|\bm R\|.
  \label{eq:gravity-model}
\end{equation}
where $\mu$ is Earth's gravitational parameter, $R_\oplus$ is the equatorial
radius, and $J_2$ is the second zonal coefficient. When $J_2$ is disabled,
the second term in~\eqref{eq:gravity-model} is omitted. The feedback
acceleration is not a physical force applied to every body; it only chooses a
convenient moving origin. With total modeled mass $m_{\rm tot}=\sum_i m_i$,
\begin{equation}
  \tilde{\bm a}_{\rm fb} =
  \frac{1}{m_{\rm tot}}\sum_i \bm F_i^{\rm ng,trans},\qquad
  \bm a_{\rm fb}=10^{-3}\tilde{\bm a}_{\rm fb},
  \label{eq:feedback}
\end{equation}
where $\bm F_i^{\rm ng,trans}$ contains drag, solar radiation pressure, and
thruster forces in SI units.

\subsection{Body Wrenches}
\label{subsec:body-wrenches}

For each \mj{} body, \mjorbit{} computes an external force in the
chief-centered world frame and hands it to \mj{} as a passive wrench. In SI
units, the translational component is
\begin{equation}
  \bm F_i =
    m_i\,10^3\!\left[\bm g(\bm R_i)-\bm g(\bm R)\right]
    -m_i\tilde{\bm a}_{\rm fb}
    +\bm F_i^{\rm ng}.
  \label{eq:body-force}
\end{equation}
The first term is differential gravity, the second compensates for the
non-gravitational acceleration of the chosen origin, and
$\bm F_i^{\rm ng}$ contains body-level environmental and actuator forces.
No Coriolis, centrifugal, or Euler-frame terms appear because the world axes
are inertially aligned.

Surface drag and solar radiation pressure are evaluated at each surface
center of pressure. For a surface normal $\hat{\bm n}_k^w$, area $A_k$,
drag coefficient $C_d^k$, reflectivity coefficient $C_r^k$, and relative
atmospheric velocity $\bm v_{ik}^{\rm rel}$ in \si{m/s},
\begin{align}
  \bm F_{ik}^{\rm drag}
   &= -\frac{1}{2}\rho(\bm R) C_d^k A_k
      \max(0,\hat{\bm n}_k^w\!\cdot\hat{\bm v}_{ik}^{\rm rel})
      \|\bm v_{ik}^{\rm rel}\|\,\bm v_{ik}^{\rm rel},\\
  \bm F_{ik}^{\rm SRP}
   &= -\eta P_\odot C_r^k A_k
      \max(0,\hat{\bm n}_k^w\!\cdot\hat{\bm s})\,\hat{\bm s}.
  \label{eq:surface-forces}
\end{align}
Here $\rho$ is atmospheric density, $\eta$ is the eclipse factor,
$P_\odot=\SI{4.56e-6}{N/m^2}$, and $\hat{\bm s}$ is the Sun direction.
The surface moment arm supplies the corresponding torque about the body COM.
Atmospheric density, magnetic field, Sun direction, and eclipse are cached at
the chief state; body position, velocity, attitude, and surface moment arms
are evaluated per body.

The rotational wrench includes finite-body gravity-gradient torque,
magnetic torques, actuator gyroscopic torques, and force-offset torques:
\begin{equation}
  \bm\tau_i =
  \bm\tau_i^{gg}
  + \bm\tau_i^{\rm surf}
  + \bm\tau_i^{\rm mag}
  + \bm\tau_i^{\rm act}
  + \bm\tau_i^{\rm thr}.
  \label{eq:body-torque-sum}
\end{equation}
The gravity-gradient term is
\begin{equation}
  \bm\tau_i^{gg}
    = \frac{3\mu}{\|\bm R_i\|^3}\,
      \hat{\bm R}_i \times \left(J_i^w\hat{\bm R}_i\right),
  \qquad \hat{\bm R}_i=\bm R_i/\|\bm R_i\|,
  \label{eq:gg-torque}
\end{equation}
where $J_i^w$ is the inertia tensor in world axes. Residual dipoles and
magnetorquers use $\bm\tau=\bm m\times\bm B$ in the body frame and rotate the
result to world coordinates. Reaction wheels and single-gimbal CMGs apply the
standard gyroscopic reaction torques from their stored momentum states, while
thrusters apply both a directed force and its moment about the body COM.

\subsection{Multibody Coupling}
\label{subsec:multibody-coupling}

\mj{} advances the generalized multibody state $(q,\dot q)$, including
free-joint roots, articulated joints, equality constraints, and contacts. With
\mj{}'s built-in uniform gravity disabled, the constrained dynamics can be
written
\begin{equation}
  M(q)\ddot q+\bm c(q,\dot q)
  =\bm\tau^{\rm passive}+\bm\tau^{\rm ctrl}+J_c^\top\bm\lambda,
  \label{eq:mj-dynamics}
\end{equation}
where $\bm\lambda$ denotes the contact and constraint impulses solved by
\mj{}. \mjorbit{} injects the orbit/environment effects through
\texttt{qfrc\_passive} using \texttt{mj\_applyFT}. Equivalently,
\begin{equation}
  \bm\tau^{\rm passive}
  =\sum_i \left(J_i^\top\bm F_i+J_i^{R\top}\bm\tau_i\right),
  \label{eq:passive-projection}
\end{equation}
with $J_i$ and $J_i^R$ the translational and rotational body Jacobians.
This projection makes the same per-body wrench formulation valid for free
bodies, articulated spacecraft, and contact-coupled systems. Contact impulses
are internal to the \mj{} solve and are not included in the chief feedback
acceleration in~\eqref{eq:feedback}.

\subsection{Step Procedure}
\label{subsec:step-procedure}

The coupling is implemented as a \mj{} plugin attached during model
compilation. A call to \texttt{mjo\_step} executes four stages:
\begin{enumerate}\setlength{\itemsep}{0pt}
  \item update environment caches from the current chief state;
  \item assemble body forces and torques from
    \eqref{eq:body-force}--\eqref{eq:body-torque-sum} into
    \texttt{qfrc\_passive};
  \item let \mj{} solve and integrate the constrained multibody dynamics;
  \item advance stored actuator states and propagate the chief state using
    \eqref{eq:chief-dynamics}, then refresh post-step caches.
\end{enumerate}
For comparisons with external simulators, \mj{} states are converted back to
absolute ECI coordinates using~\eqref{eq:absolute-state}; derived LVLH helper
transforms are used only for initialization, analysis, and visualization.

\end{document}
